% ID: SEL-D0080-OB-E02
% Tanda S080.
\textbf{a)} El dominio es $\mathbb R$. La continuidad en la unión exige $4a-4=1-b$. Como $x=1$ está dentro del tramo polinómico, el mínimo exige que se anule su derivada: $2a\cdot1-2=0$. De aquí $a=1$ y $b=1$. Con estos valores, el primer tramo es $(x-1)^2-1$, que tiene mínimo en $1$, y el segundo es creciente a partir del valor 0 en la unión. Queda confirmado el mínimo, con $\boxed{a=1,\ b=1}$.

\textbf{b)} Para los nuevos valores, el tramo izquierdo es $1{,}5x^2-2x$ y se incluye $x=2$. Tiene vértice $(2/3,-2/3)$, corta el eje horizontal en $0$ y $4/3$, y termina en el punto cerrado $(2,2)$. El tramo derecho es la recta $y=x/2-1/2$ para $x>2$, que comienza con un punto abierto en $(2,1/2)$. Hay un salto en $2$.
\begin{center}
\setlength{\unitlength}{1mm}
\begin{picture}(120,78)(-20,-28)
\put(-15,0){\vector(1,0){110}}\put(0,-23){\vector(0,1){67}}
\put(97,0){$x$}\put(1,44){$f$}
\qbezier(-10,20.625)(15,-45)(40,30)
\qbezier(40,7.5)(65,16.875)(90,26.25)
\put(40,30){\circle*{1.8}}
\put(40,7.5){\color{white}\circle*{2.4}}
\put(40,7.5){\circle{2.4}}
\put(42,31){\scriptsize $(2,2)$}
\put(42,3){\scriptsize $(2,1/2)$}
\put(7,-17){\scriptsize $(2/3,-2/3)$}
\end{picture}
\end{center}
